\documentclass[11pt, a4paper]{article}

% --- UNIVERSAL PREAMBLE BLOCK ---
\usepackage[a4paper, top=2.5cm, bottom=2.5cm, left=2cm, right=2cm]{geometry}
\usepackage{fontspec}

\usepackage[english, bidi=basic, provide=*]{babel}

\babelprovide[import]{english}

% Set default/Latin font to Sans Serif in the main (rm) slot
\babelfont{rm}{Noto Sans}

% --- ADDITIONAL PACKAGES ---
\usepackage{amssymb}    % For checkboxes \square
\usepackage{booktabs}   % For professional tables
\usepackage{longtable}  % For multi-page tracking table
\usepackage{enumitem}   % For customized lists
\usepackage{xcolor}     % For colored highlights
\usepackage[colorlinks=true, linkcolor=blue]{hyperref} % For internal links

% --- CUSTOM COMMANDS ---
\newcommand{\project}[1]{\textbf{\textcolor{blue}{Hands-on Project:}} #1}
\newcommand{\usecase}[1]{\textbf{\textcolor{teal}{Real-World Architecture:}} #1}
\newcommand{\capstone}[2]{
    \vspace{3mm}
    \noindent\fbox{
        \parbox{0.97\textwidth}{
            \textbf{\textcolor{red}{CAPSTONE PROJECT - #1:}} #2
        }
    }
    \vspace{3mm}
}
\newcommand{\dayitem}[2]{\item \textbf{#1:} #2}

\title{\Huge \textbf{Snowflake Master Data Engineer}\\ \Large 6-Month Exhaustive Architecture, Snowpark, and DataOps Program}
\author{Multi-Cluster Data Warehousing, Streaming, Governance, and AI Integration}
\date{}

\begin{document}

\maketitle

\begin{abstract}
\noindent This 24-week program is designed to transform you into a Master Snowflake Data Engineer. It covers Snowflake's core architecture, micro-partitioning, continuous ingestion with Snowpipe, CDC streams, serverless tasks, Dynamic Tables, Snowpark Python development, data governance, Snowflake Cortex AI, Hybrid Tables (Unistore), and production-grade CI/CD workflows with dbt and Terraform. A tracking timetable is included at the end.
\end{abstract}

\tableofcontents
\newpage

% ==========================================
% MILESTONE 1
% ==========================================
\section{Milestone 1: Architecture, Warehouses \& Storage Engine (Month 1)}
\textit{Objective: Master the 3-layer architecture, micro-partitions, virtual warehouses, and zero-copy cloning.}

\subsection*{Week 1: Multi-Cluster Shared Data Architecture \& Storage}
\begin{itemize}[label=-]
    \dayitem{Monday}{Snowflake's 3-Layer Architecture: Database Storage, Query Processing (Compute), and Cloud Services.}
    \dayitem{Tuesday}{Micro-partitions: Immutable columnar storage, metadata tracking, and clustering keys.}
    \dayitem{Wednesday}{Data Types, Storage Optimization, and Automatic Clustering mechanics.}
    \dayitem{Thursday}{\project{Deploy a Snowflake database, create a 10-million-row synthetic table, inspect system metadata via \texttt{SYSTEM\$CLUSTERING\_INFORMATION}, and define a custom Clustering Key.}}
    \dayitem{Friday}{\usecase{Design an enterprise storage structure for a global telecommunications company processing 5TB of CDR (Call Detail Record) logs daily, balancing auto-clustering costs against query speed.}}
\end{itemize}

\subsection*{Week 2: Virtual Warehouses \& Compute Scaling}
\begin{itemize}[label=-]
    \dayitem{Monday}{Virtual Warehouse sizing (XS to 6XL), credit consumption, and local SSD caching.}
    \dayitem{Tuesday}{Auto-suspend, Auto-resume, and Warehouse State management.}
    \dayitem{Wednesday}{Multi-Cluster Warehouses: Standard vs. Economy scaling policies, Min/Max clusters.}
    \dayitem{Thursday}{\project{Configure a Multi-Cluster Warehouse with auto-scaling (Min: 1, Max: 5). Write a Python script using the Snowflake Connector to simulate 100 concurrent queries and observe cluster scaling in snowsight.}}
    \dayitem{Friday}{\usecase{Architect a compute isolation strategy for an enterprise where BI dashboards, ETL batch processes, and ad-hoc data science queries run concurrently without compute contention.}}
\end{itemize}

\subsection*{Week 3: Tables, Views, Materialized Views \& Search Optimization}
\begin{itemize}[label=-]
    \dayitem{Monday}{Table Types: Permanent, Transient, and Temporary tables.}
    \dayitem{Tuesday}{Standard Views, Secure Views, and Materialized Views (maintenance overhead vs. speed).}
    \dayitem{Wednesday}{Search Optimization Service (SOS): Point-lookup acceleration on unstructured/structured data.}
    \dayitem{Thursday}{\project{Create Transient tables for intermediate ETL staging, build a Materialized View for daily aggregations, and enable Search Optimization on a high-cardinality UUID column.}}
    \dayitem{Friday}{\usecase{Evaluate cost vs. performance trade-offs between Automatic Clustering, Materialized Views, and Search Optimization for an e-commerce platform's analytics layer.}}
\end{itemize}

\subsection*{Week 4: Zero-Copy Cloning, Time Travel \& Fail-Safe}
\begin{itemize}[label=-]
    \dayitem{Monday}{Time Travel mechanics: Querying historical data (\texttt{AT} | \texttt{BEFORE}) up to 90 days.}
    \dayitem{Tuesday}{Fail-safe storage: Non-configurable 7-day disaster recovery tier.}
    \dayitem{Wednesday}{Zero-Copy Cloning: Metadata-only duplication of databases, schemas, and tables.}
    \dayitem{Thursday}{\project{Simulate an accidental \texttt{DROP TABLE}. Restore it using \texttt{UNDROP}. Create a Zero-Copy Clone of a Production database into a QA environment, modify data, and prove zero storage duplication.}}
    \dayitem{Friday}{\usecase{Architect a zero-cost instant staging and testing environment workflow for a DevOps pipeline using Zero-Copy Cloning and ephemeral transient schemas.}}
\end{itemize}

\capstone{Milestone 1}{Multi-Tier E-Commerce Data Backend \& Disaster Recovery System. Design an enterprise Snowflake architecture. Set up isolated multi-cluster warehouses for ETL and BI. Load raw e-commerce data into permanent tables, configure transient staging layers, set up custom clustering keys, implement Search Optimization on user IDs, and build an automated Time Travel backup testing script.}

% ==========================================
% MILESTONE 2
% ==========================================
\section{Milestone 2: Data Ingestion, Streams \& Tasks (ELT Engine) (Month 2)}
\textit{Objective: Master staging, continuous loading with Snowpipe, CDC streams, and task orchestration.}

\subsection*{Week 5: Data Staging \& Storage Integrations}
\begin{itemize}[label=-]
    \dayitem{Monday}{Internal Stages (User, Table, Named) vs. External Stages (S3, GCS, Azure Blob).}
    \dayitem{Tuesday}{Secure authentication via Storage Integrations (IAM roles without hardcoded keys).}
    \dayitem{Wednesday}{The \texttt{COPY INTO} command: File formats (CSV, JSON, Parquet, Avro, ORC) and transformation on load.}
    \dayitem{Thursday}{\project{Create an AWS IAM Role, set up a Snowflake Storage Integration, create an External Stage pointing to an S3 bucket, and execute a \texttt{COPY INTO} command with pattern matching.}}
    \dayitem{Friday}{\usecase{Design an ingestion pipeline for a logistics company receiving millions of compressed Parquet files per hour from IoT fleet sensors across AWS and Azure.}}
\end{itemize}

\subsection*{Week 6: Continuous Ingestion (Snowpipe \& Snowpipe Streaming)}
\begin{itemize}[label=-]
    \dayitem{Monday}{Snowpipe Architecture: Event-driven auto-ingestion via SQS / Event Grid notifications.}
    \dayitem{Tuesday}{Snowpipe REST API for programmatic ingestion calls.}
    \dayitem{Wednesday}{Snowpipe Streaming API: Direct row-set streaming into Snowflake tables (sub-second latency).}
    \dayitem{Thursday}{\project{Configure an S3 bucket event notification to trigger a Snowpipe automatically whenever a new CSV arrives. Verify continuous ingestion using \texttt{SYSTEM\$PIPE\_STATUS}.}}
    \dayitem{Friday}{\usecase{Architect a real-time fraud detection ingestion pipeline processing credit card swipes using Snowpipe Streaming vs. traditional micro-batching.}}
\end{itemize}

\subsection*{Week 7: Change Data Capture (CDC) with Streams}
\begin{itemize}[label=-]
    \dayitem{Monday}{Snowflake Streams: Standard, Append-Only, and Insert-Only stream types.}
    \dayitem{Tuesday}{Understanding Stream metadata columns (\texttt{METADATA\$ACTION}, \texttt{METADATA\$ISUPDATE}, \texttt{METADATA\$ROW\_ID}).}
    \dayitem{Wednesday}{Consuming streams in \texttt{MERGE} statements for incremental updates.}
    \dayitem{Thursday}{\project{Create a primary target table and a Standard Stream over a landing table. Perform \texttt{INSERT}, \texttt{UPDATE}, and \texttt{DELETE} operations on the landing table, then write a \texttt{MERGE} statement to sync changes to the target.}}
    \dayitem{Friday}{\usecase{Design an incremental CDC pipeline synchronizing an operational MySQL database into Snowflake via AWS DMS and Snowflake Streams.}}
\end{itemize}

\subsection*{Week 8: Orchestration with Snowflake Tasks}
\begin{itemize}[label=-]
    \dayitem{Monday}{Snowflake Tasks: Scheduling SQL, CRON expressions, and execution parameters.}
    \dayitem{Tuesday}{Task Trees (DAGs): Master and Child tasks using \texttt{AFTER} clauses.}
    \dayitem{Wednesday}{Serverless Tasks and Task Error Handling / Notification Integration (AWS SNS / Slack).}
    \dayitem{Thursday}{\project{Build a multi-step DAG using Tasks: Task 1 checks if a Stream has data (\texttt{SYSTEM\$STREAM\_HAS\_DATA}), Task 2 consumes the Stream via \texttt{MERGE}, and Task 3 truncates the staging environment.}}
    \dayitem{Friday}{\usecase{Architect an automated, serverless, self-healing ELT orchestration engine handling 50 dependent business data transformations without external orchestrators.}}
\end{itemize}

\capstone{Milestone 2}{Real-Time CDC \& Automated ELT Pipeline. Build a continuous ingestion engine. Configure an External Stage on AWS S3 with Snowpipe auto-ingestion into a raw JSON landing table. Attach a Standard Stream to track changes, and orchestrate a DAG using Serverless Tasks to parse JSON, execute incremental \texttt{MERGE} statements into a Gold schema, and send error alerts on failure.}

% ==========================================
% MILESTONE 3
% ==========================================
\section{Milestone 3: Advanced SQL, Semi-Structured Data \& Dynamic Tables (Month 3)}
\textit{Objective: Master VARIANT processing, query profiling, Dynamic Tables, and advanced analytics SQL.}

\subsection*{Week 9: Semi-Structured Data Processing}
\begin{itemize}[label=-]
    \dayitem{Monday}{The \texttt{VARIANT}, \texttt{OBJECT}, and \texttt{ARRAY} data types in Snowflake.}
    \dayitem{Tuesday}{Dot notation, Bracket notation, and explicit casting of semi-structured fields.}
    \dayitem{Wednesday}{\texttt{FLATTEN} and \texttt{LATERAL FLATTEN}: Unpacking nested arrays and objects into tabular rows.}
    \dayitem{Thursday}{\project{Load raw deeply nested JSON clickstream logs into a \texttt{VARIANT} column. Write queries using \texttt{LATERAL FLATTEN} to extract array elements and build a clean dimensional table.}}
    \dayitem{Friday}{\usecase{Architect a flexible schema-on-read ingestion pipeline for a SaaS company whose mobile apps emit changing JSON payload schemas without breaking production downstream models.}}
\end{itemize}

\subsection*{Week 10: Query Tuning, Performance \& Profiling}
\begin{itemize}[label=-]
    \dayitem{Monday}{Reading the Query Profile in Snowsight: Identifying bottlenecks, pruning, and scan metrics.}
    \dayitem{Tuesday}{Spilling to Local Storage vs. Spilling to Remote Storage (and how to fix it).}
    \dayitem{Wednesday}{Join optimization: Cartesians, build side vs. probe side, and search keys.}
    \dayitem{Thursday}{\project{Identify a slow query. Use the Query Profile to diagnose disk spilling. Refactor the query using CTEs, eliminate Cartesian joins, resize the warehouse, and verify execution speedup.}}
    \dayitem{Friday}{\usecase{Perform an enterprise cost and performance audit on a query pipeline consuming \$20,000/month, optimizing micro-partition pruning and warehouse allocation.}}
\end{itemize}

\subsection*{Week 11: Dynamic Tables (Declarative Data Pipelines)}
\begin{itemize}[label=-]
    \dayitem{Monday}{Dynamic Tables vs. Materialized Views vs. Tasks + Streams.}
    \dayitem{Tuesday}{Target Lag (\texttt{TARGET\_LAG = '5 minutes'}) and automated refresh mechanics.}
    \dayitem{Wednesday}{Incremental Refresh vs. Full Refresh triggers and dependency DAGs.}
    \dayitem{Thursday}{\project{Build a 3-tier pipeline (Bronze $\rightarrow$ Silver $\rightarrow$ Gold) using purely Dynamic Tables. Configure different target lags for each layer and observe automated dependency refresh.}}
    \dayitem{Friday}{\usecase{Re-architect a complex 200-task Airflow DAG into a declarative Snowflake Dynamic Tables network, reducing operational maintenance and pipeline latency.}}
\end{itemize}

\subsection*{Week 12: Advanced Analytics SQL}
\begin{itemize}[label=-]
    \dayitem{Monday}{Window Functions: \texttt{PARTITION BY}, \texttt{ORDER BY}, and Frame specifications (\texttt{ROWS} / \texttt{RANGE}).}
    \dayitem{Tuesday}{The \texttt{QUALIFY} clause: Filtering window function outputs directly without subqueries.}
    \dayitem{Wednesday}{Grouping Sets, \texttt{CUBE}, \texttt{ROLLUP}, and Pivoting / Unpivoting tables.}
    \dayitem{Thursday}{\project{Write a financial analytics query calculating a 30-day moving average and cumulative sum of customer transactions, using \texttt{QUALIFY} to retain only top-ranked records.}}
    \dayitem{Friday}{\usecase{Build an analytical reporting engine calculating customer retention cohorts and month-over-month churn across 50 million subscription billing records.}}
\end{itemize}

\capstone{Milestone 3}{Semi-Structured Telemetry Lakehouse \& Declarative Pipeline. Ingest 10GB of nested JSON IoT telemetry data. Flatten the JSON payloads into tabular Bronze Dynamic Tables. Build Silver and Gold Dynamic Tables with incremental target lags, using advanced SQL window functions and \texttt{QUALIFY} clauses to compute real-time anomaly metrics.}

\newpage
% ==========================================
% MILESTONE 4
% ==========================================
\section{Milestone 4: Snowpark Python, UDFs \& Enterprise Data Engineering (Month 4)}
\textit{Objective: Master Snowpark Python DataFrames, UDFs, Stored Procedures, and API integrations.}

\subsection*{Week 13: Snowpark Python Fundamentals}
\begin{itemize}[label=-]
    \dayitem{Monday}{Snowpark Architecture: Client-side Lazy Evaluation vs. Server-side Pushdown execution.}
    \dayitem{Tuesday}{\texttt{Snowpark Session} creation, DataFrame operations, and transformations.}
    \dayitem{Wednesday}{Translating SQL to Snowpark DataFrames: \texttt{select}, \texttt{filter}, \texttt{join}, and \texttt{group\_by}.}
    \dayitem{Thursday}{\project{Set up a local Python environment. Create a Snowpark Session, load data from a Snowflake table, build complex filter/join transformations using Snowpark DataFrames, and execute \texttt{collect()}.}}
    \dayitem{Friday}{\usecase{Design a PySpark to Snowpark migration roadmap for a enterprise data team, translating Spark DataFrames to native Snowpark Python without data extraction.}}
\end{itemize}

\subsection*{Week 14: User-Defined Functions (UDFs) \& Vectorized UDFs}
\begin{itemize}[label=-]
    \dayitem{Monday}{Scalar UDFs in Python and SQL.}
    \dayitem{Tuesday}{Vectorized UDFs using Python Type Hints (\texttt{pandas\_udf}) and PyArrow for high performance.}
    \dayitem{Wednesday}{User-Defined Table Functions (UDTFs) returning tabular outputs.}
    \dayitem{Thursday}{\project{Write a Vectorized Python UDF using Pandas to compute text sentiment scores over a million customer review rows directly inside Snowflake's compute engine.}}
    \dayitem{Friday}{\usecase{Architect a high-performance geospatial processing pipeline utilizing custom Python UDFs to perform spatial distance calculations over taxi trip data.}}
\end{itemize}

\subsection*{Week 15: Stored Procedures in Python \& Anaconda Packages}
\begin{itemize}[label=-]
    \dayitem{Monday}{UDFs vs. Stored Procedures (Data transformation vs. Administrative orchestration).}
    \dayitem{Tuesday}{Writing Python Stored Procedures using \texttt{snowflake-snowpark-python}.}
    \dayitem{Wednesday}{Leveraging third-party Python packages via Anaconda Integration (scikit-learn, numpy, pandas).}
    \dayitem{Thursday}{\project{Build a Python Stored Procedure that accepts a table name as input, uses \texttt{scikit-learn} to train a linear regression model on the data, and writes prediction results back to a target table.}}
    \dayitem{Friday}{\usecase{Design an automated ML feature engineering pipeline running entirely inside Snowflake, eliminating the need to export training data to external compute clusters.}}
\end{itemize}

\subsection*{Week 16: External Functions \& API Integrations}
\begin{itemize}[label=-]
    \dayitem{Monday}{External Functions Architecture: API Integrations + AWS API Gateway / Azure API Management.}
    \dayitem{Tuesday}{Invoking external REST APIs safely from Snowflake SQL queries.}
    \dayitem{Wednesday}{Handling API rate limits, batching, and authentication headers.}
    \dayitem{Thursday}{\project{Configure an External Function that calls an AWS Lambda endpoint to geocode addresses stored in a Snowflake table in real-time.}}
    \dayitem{Friday}{\usecase{Architect a real-time address verification and currency conversion pipeline integrating third-party SaaS APIs with Snowflake Data Engineering workflows.}}
\end{itemize}

\capstone{Milestone 4}{End-to-End Snowpark Python Engineering Pipeline. Build a complete Python-driven data pipeline. Use Snowpark Python to read staging tables, execute vectorized Pandas UDFs for text enrichment, trigger an External Function to validate postal addresses via external APIs, and wrap the entire process in a Snowpark Stored Procedure.}

% ==========================================
% MILESTONE 5
% ==========================================
\section{Milestone 5: Security, Governance, Data Sharing \& Clean Rooms (Month 5)}
\textit{Objective: Master RBAC, Column/Row level security, Data Sharing, and Snowflake Clean Rooms.}

\subsection*{Week 17: Role-Based Access Control (RBAC) Deep Dive}
\begin{itemize}[label=-]
    \dayitem{Monday}{System-defined roles (\texttt{ACCOUNTADMIN}, \texttt{SECURITYADMIN}, \texttt{SYSADMIN}, \texttt{USERADMIN}) and best practices.}
    \dayitem{Tuesday}{Custom Role Hierarchies: Functional Roles vs. Access Roles.}
    \dayitem{Wednesday}{Secondary Roles and Session Privileges.}
    \dayitem{Thursday}{\project{Design and execute a production RBAC schema: Create Access Roles for \texttt{READ\_GOLD} and \texttt{WRITE\_SILVER}, assign them to Functional Roles (\texttt{ANALYST}, \texttt{ENGINEER}), and grant to Users.}}
    \dayitem{Friday}{\usecase{Architect an enterprise security hierarchy for 500 users across 10 business units, strictly enforcing least-privilege access and separating administrative duties.}}
\end{itemize}

\subsection*{Week 18: Column \& Row Level Security, Data Tagging}
\begin{itemize}[label=-]
    \dayitem{Monday}{Object Tagging: Assigning metadata tags to databases, schemas, tables, and columns.}
    \dayitem{Tuesday}{Dynamic Data Masking Policies: Obfuscating PII data based on user roles.}
    \dayitem{Wednesday}{Row Access Policies: Restricting row-level visibility based on user context (\texttt{CURRENT\_ROLE()}).}
    \dayitem{Thursday}{\project{Create a Tag-Based Masking Policy that automatically masks SSN columns for non-HR roles. Build a Row Access Policy restricting sales reps to seeing only their region's rows.}}
    \dayitem{Friday}{\usecase{Design a HIPAA/GDPR compliance framework in Snowflake that automatically classifies sensitive PII data and enforces dynamic masking across all analytical schemas.}}
\end{itemize}

\subsection*{Week 19: Snowflake Data Sharing \& Marketplace}
\begin{itemize}[label=-]
    \dayitem{Monday}{Snowflake Secure Data Sharing Architecture (Zero-copy live data access without ETL).}
    \dayitem{Tuesday}{Direct Shares vs. Reader Accounts (sharing data with non-Snowflake customers).}
    \dayitem{Wednesday}{Snowflake Marketplace: Monetizing data products and consuming third-party listings.}
    \dayitem{Thursday}{\project{Create a Secure Share, add secure views to the share, grant access to a simulated consumer account, and verify zero data copying.}}
    \dayitem{Friday}{\usecase{Architect a B2B data distribution platform for a financial data provider, serving live stock market feeds to 100 enterprise clients using Reader Accounts.}}
\end{itemize}

\subsection*{Week 20: Snowflake Data Clean Rooms \& Native Apps}
\begin{itemize}[label=-]
    \dayitem{Monday}{Data Clean Room concepts: Multi-party privacy-preserving analytics without raw data exposure.}
    \dayitem{Tuesday}{Differential Privacy and cryptographic hashing in Clean Rooms.}
    \dayitem{Wednesday}{Introduction to the Snowflake Native App Framework.}
    \dayitem{Thursday}{\project{Build a two-party Data Clean Room scenario: Allow an advertiser and a publisher to intersect customer lists and compute overlap metrics without exposing email addresses.}}
    \dayitem{Friday}{\usecase{Design a healthcare research Data Clean Room allowing hospital networks and pharmaceutical companies to run joint ML studies without sharing raw patient records.}}
\end{itemize}

\capstone{Milestone 5}{Governed Multi-Tenant Data Platform \& Secure Share. Build a production security perimeter. Implement an Access/Functional RBAC role hierarchy, create Tag-Based Dynamic Data Masking for PII, enforce Row Access Policies for regional multi-tenancy, and build a Secure Data Share to expose anonymized aggregated Gold datasets to external Reader Accounts.}

\newpage
% ==========================================
% MILESTONE 6
% ==========================================
\section{Milestone 6: Cortex AI, Unistore, DataOps \& Certification (Month 6)}
\textit{Objective: Master Cortex AI, Hybrid Tables, Terraform/dbt DataOps workflows, and SnowPro certification.}

\subsection*{Week 21: Unistore \& Hybrid Tables (Hybrid OLTP + OLAP)}
\begin{itemize}[label=-]
    \dayitem{Monday}{Unistore Architecture: Combining transactional (OLTP) and analytical (OLAP) workloads.}
    \dayitem{Tuesday}{Hybrid Tables: Primary Keys, Foreign Keys, Unique Constraints, and Indexes.}
    \dayitem{Wednesday}{Row-level locking and single-digit millisecond point lookups in Hybrid Tables.}
    \dayitem{Thursday}{\project{Create a Hybrid Table with Primary Key constraints and secondary indexes. Execute high-frequency single-row \texttt{INSERT} and \texttt{UPDATE} queries, then join it directly with an OLAP analytical table.}}
    \dayitem{Friday}{\usecase{Architect an e-commerce order management system where real-time order writes hit Hybrid Tables and are immediately available for analytical dashboards without CDC delays.}}
\end{itemize}

\subsection*{Week 22: Snowflake Cortex AI \& Vector Search}
\begin{itemize}[label=-]
    \dayitem{Monday}{Snowflake Cortex AI: Serverless LLM functions (\texttt{CORTEX.COMPLETE}, \texttt{CORTEX.EXTRACT\_ANSWER}).}
    \dayitem{Tuesday}{\texttt{CORTEX.EMBED\_TEXT\_768} and generating vector embeddings natively inside SQL.}
    \dayitem{Wednesday}{Vector Data Type and Vector Similarity Search for RAG (Retrieval-Augmented Generation).}
    \dayitem{Thursday}{\project{Build a native RAG pipeline inside Snowflake: Extract text from PDFs using \texttt{PARSE\_DOCUMENT}, generate vector embeddings with Cortex AI, store them in a Vector column, and search using vector distance functions.}}
    \dayitem{Friday}{\usecase{Design an enterprise document intelligence pipeline that automatically parses incoming PDF invoices, extracts line items via Cortex AI, and indexes them for vector search.}}
\end{itemize}

\subsection*{Week 23: Production DataOps (dbt, Terraform \& Schemachange)}
\begin{itemize}[label=-]
    \dayitem{Monday}{Managing Snowflake infrastructure as code using Terraform.}
    \dayitem{Tuesday}{Database Migration Management with \texttt{schemachange} / Flyway.}
    \dayitem{Wednesday}{Integrating Snowflake with dbt (data build tool): Models, tests, and documentation.}
    \dayitem{Thursday}{\project{Set up a GitHub repository with dbt and \texttt{schemachange}. Build a CI/CD pipeline using GitHub Actions that automatically deploys Snowflake table changes and runs dbt models on pull requests.}}
    \dayitem{Friday}{\usecase{Architect a multi-environment (Dev, QA, Prod) DataOps deployment framework ensuring automated testing, zero-downtime schema migrations, and Git-driven approval workflows.}}
\end{itemize}

\subsection*{Week 24: SnowPro Advanced Data Engineer Exam Preparation}
\begin{itemize}[label=-]
    \dayitem{Monday}{Reviewing SnowPro Advanced Data Engineer Domain 1: Data Movement (30\%).}
    \dayitem{Tuesday}{Reviewing Domain 2: Performance Tuning (30\%) \& Domain 3: Storage and Protection (20\%).}
    \dayitem{Wednesday}{Reviewing Domain 4: Security and Governance (20\%) \& Case Studies.}
    \dayitem{Thursday}{\project{Take a full-length 65-question practice exam under timed conditions. Review missed questions, focusing on obscure syntax, clustering depth math, and warehouse scaling edge cases.}}
    \dayitem{Friday}{Final Review and Certification Readiness Verification.}
\end{itemize}

\capstone{Milestone 6}{Production-Grade Autonomous Data Mesh Engine. Build an enterprise-ready Data Mesh architecture. Deploy infrastructure via Terraform. Build a Dynamic Table transformation layer managed by dbt and GitHub Actions CI/CD. Integrate Hybrid Tables for transaction processing, implement Cortex AI for LLM text enrichment, lock down the platform with RBAC and Dynamic Data Masking, and execute a dry-run SnowPro practice evaluation.}

\newpage
% ==========================================
% TRACKING TIMETABLE
% ==========================================
\section{6-Month Progress Tracking Timetable}

\begin{center}
\renewcommand{\arraystretch}{1.5}
\begin{longtable}{@{} c l p{7cm} c @{}}
\toprule
\textbf{Week} & \textbf{Primary Focus} & \textbf{Key Project / Capstone} & \textbf{Status} \\ 
\midrule
\endfirsthead

\toprule
\textbf{Week} & \textbf{Primary Focus} & \textbf{Key Project / Capstone} & \textbf{Status} \\ 
\midrule
\endhead

\bottomrule
\endfoot

\bottomrule
\endlastfoot

1 & Architecture \& Micro-partitions & Clustering Information Metadata Analysis & $\square$ \\
2 & Virtual Warehouses & Multi-Cluster Auto-scaling Benchmark & $\square$ \\
3 & Tables, Views \& Search Optimization & Search Optimization Service Setup & $\square$ \\
4 & Time Travel \& Zero-Copy Clone & \textbf{Cap 1: Multi-Tier E-Commerce Backend} & $\square$ \\
\midrule
5 & Data Staging \& Integration & S3 External Stage \& Storage Integration & $\square$ \\
6 & Snowpipe \& Continuous Ingestion & SQS Event-Driven Pipe Auto-Ingestion & $\square$ \\
7 & Change Data Capture (Streams) & Standard Stream \& MERGE Synchronization & $\square$ \\
8 & Snowflake Tasks & \textbf{Cap 2: Real-Time CDC \& ELT Pipeline} & $\square$ \\
\midrule
9 & Semi-Structured Data & VARIANT JSON Flattening \& Parsing & $\square$ \\
10 & Query Profile \& Tuning & Spilling Diagnosis \& Query Refactoring & $\square$ \\
11 & Dynamic Tables & Bronze-Silver-Gold Target Lag Engine & $\square$ \\
12 & Advanced Analytics SQL & \textbf{Cap 3: Telemetry Lakehouse} & $\square$ \\
\midrule
13 & Snowpark Python Fundamentals & Snowpark Session \& Pushdown Execution & $\square$ \\
14 & Vectorized Python UDFs & Pandas Vectorized Sentiment UDF & $\square$ \\
15 & Python Stored Procedures & In-Database ML Training Procedure & $\square$ \\
16 & External Functions & \textbf{Cap 4: Snowpark Python Engineering} & $\square$ \\
\midrule
17 & Role-Based Access Control & Functional \& Access Role Hierarchy & $\square$ \\
18 & Column/Row Level Governance & Tag-Based Dynamic Data Masking & $\square$ \\
19 & Secure Data Sharing & Direct Share \& Reader Account Config & $\square$ \\
20 & Data Clean Rooms & \textbf{Cap 5: Governed Multi-Tenant Platform} & $\square$ \\
\midrule
21 & Unistore \& Hybrid Tables & Hybrid Table Transaction Execution & $\square$ \\
22 & Cortex AI \& Vector Search & Native SQL RAG Pipeline with Cortex & $\square$ \\
23 & DataOps (dbt \& Terraform) & GitHub Actions CI/CD for Snowflake & $\square$ \\
24 & Certification Preparation & \textbf{Cap 6: Autonomous Data Mesh Engine} & $\square$ \\

\end{longtable}
\end{center}

\vspace{1cm}
\begin{center}
    \textit{"In Snowflake, compute is ephemeral, metadata is king, and simplicity scales."} \\
    Best of luck on your 6-month journey to becoming a Master Snowflake Data Engineer!
\end{center}

\end{document}